% ===========================================================================
\section{AEP and typical sequences}\label{sec:aep}
\warmth{0}\quad\whisper{A long random block behaves as though its likely outcomes were nearly equiprobable.}

\begin{strand}
The asymptotic equipartition property turns entropy into a geometric picture:
almost all probability lies on roughly $2^{n\Ent(X)}$ length-$n$ sequences,
and each of those sequences has probability roughly $2^{-n\Ent(X)}$.
\end{strand}

\trigger{Use typical sets when a one-symbol distribution must be converted into a statement about long blocks.}

\block{A Bernoulli preview}

\begin{example}[counting zeros and ones]
Let $X\sim\operatorname{Bernoulli}(p)$ and let $X_1,\ldots,X_n$ be i.i.d.
copies. If a realization $x^n=(x_1,\ldots,x_n)$ contains $N_1(x^n)$ ones,
then
\[
  p_{X^n}(x^n)=p^{N_1(x^n)}(1-p)^{n-N_1(x^n)}.
\]
The law of large numbers says $N_1(X^n)/n\to p$ in probability. Thus, for
most long realizations,
\[
\begin{aligned}
  -\frac1n\log p_{X^n}(X^n)
  &\approx -p\log p-(1-p)\log(1-p)\\
  &=\answerin[1.8cm]{\Ent(X)}.
\end{aligned}
\]
\studentrecall{The limit is the binary entropy $\Ent(X)=h_2(p)$.}
Equivalently, a typical realization has probability about $2^{-n\Ent(X)}$.
\end{example}

\block{Three modes of random-variable convergence}

\begin{definition}[convergence in probability]\label{def:convergence-in-probability}
Let $Y_1,Y_2,\ldots$ and $Y$ be real-valued random variables on the same
probability space. We say that $Y_n$ \keyword{converges in probability} to $Y$,
and write $Y_n\xrightarrow{\mathrm{P}}Y$, if for every $\epsilon>0$,
\[
  \PP\bigl(|Y_n-Y|>\epsilon\bigr)
  \longrightarrow \answerin[1.2cm]{0}
  \qquad\text{as }n\to\infty.
\]
\studentrecall{For every fixed tolerance, the probability of exceeding it tends to $0$.}
Thus convergence in probability controls the probability of a fixed-size error;
it does not assert that every sample path eventually stays close.
\end{definition}

\begin{definition}[mean-square convergence]\label{def:mean-square-convergence}
Suppose $Y_1,Y_2,\ldots$ and $Y$ are square-integrable. We say that $Y_n$
\keyword{converges in mean square}, or in $L^2$, to $Y$ if
\[
  \E\!\left[|Y_n-Y|^2\right]\longrightarrow0.
\]
The notation is $Y_n\xrightarrow{L^2}Y$. This asks the average squared error,
not merely the probability of a large error, to vanish.
\end{definition}

\clearpage
\begin{definition}[almost-sure convergence]\label{def:almost-sure-convergence}
We say that $Y_n$ \keyword{converges almost surely}, or with probability one,
to $Y$ if
\[
  \PP\!\left(\left\{\omega:
  \lim_{n\to\infty}Y_n(\omega)=Y(\omega)\right\}\right)=1.
\]
The notation is $Y_n\xrightarrow{\mathrm{a.s.}}Y$. Thus, outside a single
probability-zero exceptional set, the numerical sequence $Y_n(\omega)$
converges to $Y(\omega)$ along every remaining sample path.
\end{definition}

The implications needed most often are
\[
  Y_n\xrightarrow{L^2}Y
  \quad\Longrightarrow\quad
  Y_n\xrightarrow{\mathrm{P}}Y,
  \qquad
  Y_n\xrightarrow{\mathrm{a.s.}}Y
  \quad\Longrightarrow\quad
  Y_n\xrightarrow{\mathrm{P}}Y.
\]
The first follows from
$\PP(|Y_n-Y|>\epsilon)\le\E[|Y_n-Y|^2]/\epsilon^2$.
In general, mean-square and almost-sure convergence do not imply one another
without extra assumptions.

In the AEP, $Y_n=-n^{-1}\log p_{X^n}(X_1,\ldots,X_n)$ and the limit
$Y=\Ent(X)$ is constant; the theorem asserts convergence in probability.

\begin{theorem}[weak asymptotic equipartition property]\label{thm:weak-aep}
Let $X$ be a discrete random variable on a finite alphabet, and let
$X_1,X_2,\ldots$ be i.i.d. copies of $X$. Then
\[
  -\frac1n\log p_{X^n}(X_1,\ldots,X_n)
  \longrightarrow \Ent(X)
  \qquad\text{in probability}.
\]
\end{theorem}

\begin{proof}
On the support of $X$, define $Z_i=-\log p_X(X_i)$. Independence gives
\[
  -\log p_{X^n}(X_1,\ldots,X_n)
  =\sum_{i=1}^n-\log p_X(X_i)=\sum_{i=1}^n Z_i.
\]
The variables $Z_i$ are i.i.d. and
\[
  \E[Z_i]=\sum_xp_X(x)\log\frac1{p_X(x)}=\Ent(X).
\]
The weak law of large numbers applied to $n^{-1}\sum_iZ_i$ proves the claim.
\end{proof}

\begin{yourturn}
What is the one-symbol random variable to which the law of large numbers is
applied in the AEP proof?
\studentrecall{Use the surprise $Z=-\log p_X(X)$, whose mean is $\Ent(X)$.}
\ifshowanswers\else\workspace[1]\fi
\answerprose{Apply the law of large numbers to the surprise
$Z=-\log p_X(X)$. Its expectation is exactly $\Ent(X)$.}
\end{yourturn}

\block{Weak typicality: make ``about'' precise}

\begin{definition}[weakly typical set]\label{def:weak-typical-set}
For $\epsilon>0$, the length-$n$ weakly typical set is
\[
  T_{n,\epsilon}
  =\left\{x^n\in\X^n:
  \left|-\frac1n\log p_{X^n}(x^n)-\Ent(X)\right|
  \le \answerin[1.2cm]{\epsilon}\right\}.
\]
\studentrecall{The permitted per-symbol error is $\epsilon$.}
Its members are the $\epsilon$-typical sequences of length $n$.
\end{definition}

\begin{theorem}[typical-set properties]\label{thm:typical-set-properties}
For every $\epsilon>0$, all sufficiently large $n$ satisfy:
\begin{enumerate}
  \item every $x^n\in T_{n,\epsilon}$ obeys
  \[
    2^{-n(\Ent(X)+\epsilon)}
    \le p_{X^n}(x^n)
    \le 2^{-n(\Ent(X)-\epsilon)};
  \]
  \item $\PP(X^n\in T_{n,\epsilon})>1-\epsilon$;
  \item the cardinality is bounded by
  \[
    (1-\epsilon)2^{n(\Ent(X)-\epsilon)}
    <|T_{n,\epsilon}|
    \le2^{n(\Ent(X)+\epsilon)}.
  \]
\end{enumerate}
The first property follows from the definition for every $n$; only the latter
two require $n$ to be large.
\end{theorem}

\begin{proof}
The first claim is obtained by exponentiating the defining inequality. The AEP
gives
\[
  \PP(X^n\notin T_{n,\epsilon})
  =\PP\!\left(\left|-\frac1n\log p_{X^n}(X^n)-\Ent(X)\right|>\epsilon\right)
  \longrightarrow0,
\]
which proves the second claim.

For the upper cardinality bound, use the lower probability bound on each typical
sequence:
\[
  1\ge\sum_{x^n\in T_{n,\epsilon}}p_{X^n}(x^n)
  \ge |T_{n,\epsilon}|\,2^{-n(\Ent(X)+\epsilon)}.
\]
For the lower cardinality bound, use the upper probability bound together with
the mass estimate:
\[
  1-\epsilon
  <\sum_{x^n\in T_{n,\epsilon}}p_{X^n}(x^n)
  \le |T_{n,\epsilon}|\,2^{-n(\Ent(X)-\epsilon)}.
\]
Rearranging gives the stated bounds.
\end{proof}

\begin{yourturn}
Complete the AEP slogan by combining the probability and cardinality bounds:
there are about \answerin[2.6cm]{$2^{n\Ent(X)}$} typical sequences, each with
probability about \answerin[2.6cm]{$2^{-n\Ent(X)}$}.
\studentrecall{Count $\approx2^{n\Ent(X)}$; probability $\approx2^{-n\Ent(X)}$.}
\end{yourturn}

\block{The smallest set carrying almost all mass}

\begin{definition}[minimum high-probability set]\label{def:minimum-high-probability-set}
For $0<\delta<1$, let $S_m^\delta\subseteq\X^m$ be a set of
\answerin[3.4cm]{minimum cardinality} subject to
\[
  \PP(X^m\in S_m^\delta)\ge1-\delta.
\]
\studentrecall{Minimize cardinality while retaining probability at least $1-\delta$.}
This is the ``strong'' viewpoint in the handwritten notes: keep only the most
probable sequences needed to capture the prescribed mass.
\end{definition}

\begin{proposition}[exponential equivalence of typical sets]\label{prop:typical-set-equivalence}
Let $(\delta_j)$ be any positive sequence with $\delta_j\to0$. Then
\[
  \lim_{j\to\infty}\limsup_{m\to\infty}
  \left|\frac1m\log
  \frac{|S_m^{\delta_j}|}{|T_{m,\delta_j}|}\right|=0.
\]
Thus the minimum high-probability set and the weakly typical set may differ in
their individual members, but not in their exponential size as the tolerance
shrinks.
\end{proposition}

\begin{proof}
For each fixed $0<\delta<1$, the AEP and the typical-set bounds imply
\[
  \lim_{m\to\infty}\frac1m\log|S_m^\delta|=\Ent(X),
  \qquad
  \Ent(X)-\delta
  \le\liminf_m\frac1m\log|T_{m,\delta}|
  \le\limsup_m\frac1m\log|T_{m,\delta}|
  \le\Ent(X)+\delta.
\]
For the first limit, an upper bound comes from choosing a typical set; a lower
bound follows by intersecting any set of mass at least $1-\delta$ with an
$\eta$-typical set and then letting $\eta\downarrow0$. Hence the normalized
logarithm of the ratio has absolute limsup at most $\delta$. Substituting
$\delta=\delta_j$ and letting $j\to\infty$ proves the result.
\end{proof}

\recall{How do ``probability per typical sequence'' and ``number of typical sequences'' multiply back to one?}
